\documentclass[12pt,a4paper]{article}

% --- Packages ---
\usepackage[utf8]{inputenc}
\usepackage[margin=1in]{geometry}
\usepackage{graphicx}
\usepackage{enumitem}
\usepackage{tabularx}
\usepackage{setspace}
\usepackage{hyperref}
\usepackage{listings}
\usepackage{xcolor}

\setstretch{1.2} % Line spacing

% --- Code Listing Setup ---
\lstset{
    language=Java,
    basicstyle=\ttfamily\small,
    keywordstyle=\color{blue},
    commentstyle=\color{gray},
    stringstyle=\color{red},
    breaklines=true,
    breakatwhitespace=true,
    tabsize=4,
    showstringspaces=false
}

\begin{document}

% ==========================================
% COVER PAGE
% ==========================================
\begin{titlepage}
    \centering
    \vspace*{0.5cm}

    % Institution placeholder (can add logo here)
    \framebox(2cm,1cm){[INSTITUTION LOGO]}
    \vspace{1cm}

    {\Huge \textbf{Smart City Traffic Simulator}} \\
    \vspace{0.3cm}
    {\large \textit{Adaptive Signal Control with Emergency Priority}} \\
    \vspace{1.3cm}

    {\Large \textbf{Course: Object-Oriented Programming with Java (OOPJ)}} \\
    \vspace{1.2cm}

    \textbf{Submitted By:} \\
    \vspace{0.4cm}
    {\large Anupam  \quad Arjun  \quad Abhinav} \\
    \vspace{0.2cm}
    \textit{Roll Numbers: [To be filled]} \\

    \vspace{1.5cm}

    \textbf{Submitted To:} \\
    {\large Dr./Prof. [Name]} \\

    \vspace{1.5cm}

    Department of [To be filled] \\
    \vspace{0.2cm}
    [University Name, To be filled] \\
    \vspace{0.2cm}
    Semester: [To be filled] \\

    \vfill

    \textbf{Submission Date:} \\
    {\large \today}

\end{titlepage}

\newpage

% ==========================================
% TABLE OF CONTENTS
% ==========================================
\tableofcontents
\newpage

% ==========================================
% ABSTRACT
% ==========================================
\section*{Abstract}
This project presents a Smart City Traffic Simulator that models vehicle movement and signal control at a four-way urban intersection. The primary goal is to demonstrate how object-oriented design and event-driven programming can be used to solve a practical traffic-management problem. The system is implemented in Java, with an interactive desktop interface built using Java Swing. Key features include multi-lane signal coordination, adaptive lane switching when traffic is low, emergency-vehicle priority override, real-time lane statistics, simulation speed control, pause/resume support, and session log viewing. The implementation also includes persistent file logging through a dedicated logger module. From an academic perspective, the project applies abstraction, inheritance, polymorphism, interfaces, constructor/method overloading, collections, and exception handling in a single coherent application. The simulator therefore serves as both a functional prototype and a strong learning model for OOPJ coursework.

\newpage

% ==========================================
% PROBLEM STATEMENT
% ==========================================
\section{Problem Statement}
Urban traffic congestion at peak hours negatively impacts air quality, fuel consumption, and emergency response times. Most traditional traffic signal systems operate on fixed timing cycles, leading to inefficient green-light allocation. When a lane is empty, the green light is still held for the full configured duration, while waiting traffic in other lanes grows unnecessarily long. Furthermore, emergency vehicles like ambulances are often forced to wait through red signals, creating life-threatening delays.

We designed the Smart City Traffic Simulator to study how adaptive signal control and emergency preemption can improve intersection efficiency. The system uses a rule-based intelligent controller that detects when active lanes are empty and can switch early, accumulating ``saved time'' in a time bank that is then reused to extend green phases for congested lanes. This approach mirrors real-world Adaptive Traffic Signal Control (ATSC) systems deployed in major cities.

The project is implemented as a complete Java application with a real-time Swing GUI, bringing the theoretical benefits of smart traffic control into a tangible, interactive learning environment. For OOPJ coursework, it serves as an excellent vehicle to demonstrate practical OOP design: abstract vehicle hierarchies, event-driven control loops, modular decomposition, exception handling, and persistent logging all work together seamlessly.

\subsection{What Problem Does This Solve?}
\begin{itemize}
    \item \textbf{Wasted Green Time:} Empty active lanes no longer hold green indefinitely; the system switches proactively.
    \item \textbf{Improved Throughput:} Accumulating saved time creates green extensions for congested lanes, reducing average wait time.
    \item \textbf{Emergency Response:} Ambulances can trigger immediate signal override for faster passage through intersections.
    \item \textbf{Real-Time Visibility:} Per-lane vehicle counts and timers let operators (or students) understand live traffic state and controller decisions.
    \item \textbf{Practical OOP Learning:} The single-intersection model is simple enough to understand but complex enough to apply 8+ core OOP concepts.
\end{itemize}

% ==========================================
% FEATURES LIST
% ==========================================
\section{Features List}

\subsection{Traffic Control}
\begin{table}[h]
    \centering
    \begin{tabularx}{\textwidth}{|l|X|}
        \hline
        \textbf{Feature} & \textbf{Description} \\ \hline
        Fixed-Cycle Mode & 80-second cycle with equal green time per lane (North → East → South → West) \\ \hline
        Smart Adaptive Mode & Early lane switch when active lane is empty; saved time applied to congested lanes \\ \hline
        Emergency Override & Ambulance detection triggers immediate green or forced yellow transition \\ \hline
        Time-Bank Management & Accumulates leftover green time and applies it intelligently to heavy-traffic lanes \\ \hline
    \end{tabularx}
\end{table}

\subsection{Vehicle Types}
\begin{table}[h]
    \centering
    \begin{tabularx}{\textwidth}{|l|l|X|}
        \hline
        \textbf{Vehicle Type} & \textbf{Color} & \textbf{Behavior} \\ \hline
        Private Car & Blue & Standard behavior; obeys red/yellow/green signals strictly \\ \hline
        Public Bus & Yellow & Larger size; follows same signal rules as private cars \\ \hline
        Motorcycle & Green & Higher acceleration; aggressive yellow crossing (more likely to enter on yellow) \\ \hline
        Ambulance & Red & Emergency behavior; can trigger signal preemption; bypasses normal red-light rules \\ \hline
    \end{tabularx}
\end{table}

\subsection{User Interface Controls}
\begin{table}[h]
    \centering
    \begin{tabularx}{\textwidth}{|l|X|}
        \hline
        \textbf{Control} & \textbf{Function} \\ \hline
        Spawn Vehicles (per lane) & Add random vehicles to a specific lane for testing \\ \hline
        Spawn Ambulance & Add an emergency vehicle to a chosen lane \\ \hline
        Toggle Smart Mode & Enable/disable adaptive lane switching and time-bank logic \\ \hline
        Speed Slider & Adjust simulation playback speed (0.5x to 2x) \\ \hline
        Pause / Resume & Freeze or continue simulation timers \\ \hline
        Clear Traffic & Remove all vehicles and reset lane counts \\ \hline
        View Logs & Open dialog showing session timestamps and debug messages \\ \hline
        Quit App & Graceful shutdown with timer cleanup and exit message \\ \hline
    \end{tabularx}
\end{table}

\subsection{Real-Time Monitoring}
\begin{table}[h]
    \centering
    \begin{tabularx}{\textwidth}{|l|X|}
        \hline
        \textbf{Metric} & \textbf{Display and Update Rate} \\ \hline
        Per-Lane Vehicle Count & Updated every 50ms (animation cycle); shows live queue depth \\ \hline
        Current Signal State & Displays active lane, color (R/Y/G), and countdown timer \\ \hline
        Time Bank Balance & Shows accumulated saved time in smart mode \\ \hline
        Emergency Override Count & Increments each time ambulance preemption is triggered \\ \hline
        Simulation Mode & Indicator showing fixed vs. smart logic active \\ \hline
    \end{tabularx}
\end{table}

% ==========================================
% INTRODUCTION
% ==========================================
\newpage
\section{Introduction and Motivation}
Rapid urbanization has significantly increased road congestion, especially at busy intersections where fixed-time traffic logic is often unable to react to changing vehicle density. In real environments, this leads to longer waiting times, inefficient green-signal usage, and delayed emergency response.

Traditional signal systems generally follow a static cycle. While easy to manage, static systems waste green time when a lane is empty and still force heavily loaded lanes to wait for the next cycle. To study and improve this behavior in a controlled setting, this project models a smart intersection simulator with adaptive logic and emergency prioritization.

The simulator is designed as an academic software system where traffic entities (vehicles, lanes, signals, controller) are represented as collaborating objects. This makes the project suitable for demonstrating core OOP principles in a realistic domain. At the same time, the GUI and timer-driven simulation loop allow users to observe immediate cause-effect relationships between control decisions and traffic flow.

\section{Objectives}
\begin{itemize}
    \item Design and implement a modular smart-intersection simulator using Java.
    \item Apply OOP principles (abstraction, inheritance, polymorphism, interfaces, encapsulation) in a real problem context.
    \item Build an interactive GUI for user control, visual feedback, and live monitoring.
    \item Implement adaptive signal behavior to reduce wasted green time on empty lanes.
    \item Implement emergency override so ambulances receive immediate movement priority.
    \item Validate functional behavior using structured test cases and observed simulation output.
\end{itemize}

\section{System Requirements}
\subsection{Hardware Requirements}
\begin{itemize}
    \item Computer/Laptop
    \item Minimum 4 GB RAM
    \item Display: 1024x768 or higher resolution for comfortable GUI interaction
\end{itemize}

\subsection{Software Requirements}
\begin{itemize}
    \item Java JDK 8 or higher (JDK 11+ recommended)
    \item IDE: NetBeans, Eclipse, or IntelliJ IDEA
    \item Operating System: Windows, Linux, or macOS with JVM support
\end{itemize}

For packaging and cross-platform installer generation in CI/CD, newer JDK toolchains (for \texttt{jpackage}) are recommended.

% ==========================================
% OOP CONCEPTS DEMONSTRATED
% ==========================================
\newpage
\section{Object-Oriented Programming Concepts Used}

This project demonstrates all core OOPJ principles through practical application:

\subsection{Abstraction – Abstract Vehicle Base Class}
The \texttt{Vehicle} class is declared abstract, defining the contract that all vehicle types must follow while leaving implementation to subclasses:

\begin{lstlisting}
public abstract class Vehicle {
    protected int x, y, speed;
    protected Color color;
    protected String laneID;
    
    // Abstract method - must be overridden by all subclasses
    public abstract void move(String signalState);
    
    public Vehicle(int x, int y, int speed, Color color, String laneID) {
        this.x = x;
        this.y = y;
        this.speed = speed;
        this.color = color;
        this.laneID = laneID;
    }
}
\end{lstlisting}

This forces all vehicle subclasses to implement their own movement logic while inheriting shared state and utility methods.

\subsection{Inheritance – Vehicle Subclasses}
Four concrete vehicle types extend the abstract base class, each with unique movement behaviors:

\begin{lstlisting}
public class PrivateCar extends Vehicle {
    public PrivateCar(int x, int y) {
        super(x, y, 4, Color.BLUE, "");
    }
    
    @Override
    public void move(String signalState) {
        if (signalState.equals("GREEN")) {
            x += speed;
        } else if (signalState.equals("YELLOW")) {
            // Conservative: 30% chance to proceed
            if (Math.random() < 0.3) x += speed;
        }
        // RED: do not move
    }
}

public class Motorcycle extends Vehicle {
    public Motorcycle(int x, int y) {
        super(x, y, 6, Color.GREEN, "");
    }
    
    @Override
    public void move(String signalState) {
        if (signalState.equals("GREEN")) {
            x += speed;
        } else if (signalState.equals("YELLOW")) {
            // Aggressive: 70% chance to proceed on yellow
            if (Math.random() < 0.7) x += speed;
        }
    }
}
\end{lstlisting}

\subsection{Polymorphism – Same Call, Different Behavior}
The GUI calls \texttt{move()} on every vehicle without knowing their type. Each subclass overrides movement differently:

\begin{lstlisting}
// In simulation loop - works with ANY vehicle type
for (Vehicle v : vehicles) {
    String currentSignal = trafficController.getCurrentSignal(v.getLaneID());
    v.move(currentSignal);  // Polymorphic call
}
\end{lstlisting}

\subsection{Interfaces – EmergencyVehicle Behavior}
The \texttt{EmergencyVehicle} interface defines special capabilities for emergency vehicles:

\begin{lstlisting}
public interface EmergencyVehicle {
    void activateEmergency();
    void deactivateEmergency();
    boolean isEmergencyActive();
}

public class Ambulance extends Vehicle implements EmergencyVehicle {
    private boolean emergencyActive = false;
    
    @Override
    public void activateEmergency() {
        emergencyActive = true;
        trafficController.handleEmergencyOverride(this.laneID);
    }
    
    @Override
    public boolean isEmergencyActive() {
        return emergencyActive;
    }
    
    @Override
    public void move(String signalState) {
        if (emergencyActive) {
            x += 8;  // Always move at max speed
        } else {
            // Normal movement - treat like motorcycle
            super.move(signalState);
        }
    }
}
\end{lstlisting}

\subsection{Encapsulation – Private Fields and Validated Accessors}
All fields in model classes are private with public getter/setter methods:

\begin{lstlisting}
public class TrafficLight {
    private String laneID;
    private int greenDuration;
    private int countdownSeconds;
    private String state; // RED, YELLOW, GREEN
    
    public void setCountdownSeconds(int seconds) {
        if (seconds < 0 || seconds > 100) {
            throw new IllegalArgumentException("Invalid countdown");
        }
        this.countdownSeconds = seconds;
    }
    
    public int getCountdownSeconds() {
        return countdownSeconds;
    }
    
    public String getState() {
        return state;
    }
}
\end{lstlisting}

\subsection{Constructor/Method Overloading}
Vehicle constructors provide flexible initialization:

\begin{lstlisting}
// Minimal constructor
public Vehicle(int x, int y, int speed) {
    this(x, y, speed, Color.BLUE, "NORTH");
}

// Full constructor with all parameters
public Vehicle(int x, int y, int speed, Color color, String laneID) {
    this.x = x;
    this.y = y;
    this.speed = speed;
    this.color = color;
    this.laneID = laneID;
}
\end{lstlisting}

\subsection{Collections – HashMap, ArrayList}
The system uses appropriate collection types for different purposes:

\begin{lstlisting}
// HashMap: Fast O(1) lane lookup for traffic lights
private Map<String, TrafficLight> trafficLights = new HashMap<>();
trafficLights.put("NORTH", new TrafficLight("NORTH", 20));
trafficLights.put("EAST", new TrafficLight("EAST", 20));

// ArrayList: Ordered vehicle list for iteration
private List<Vehicle> vehicles = new ArrayList<>();
vehicles.add(new PrivateCar(100, 200));
vehicles.add(new Ambulance(150, 200));

// Stream operations: Filter and count vehicles in a lane
long vehicleCountInNorth = vehicles.stream()
    .filter(v -> v.getLaneID().equals("NORTH"))
    .count();
\end{lstlisting}

\subsection{Exception Handling – Custom Exceptions}
The project defines domain-specific exceptions:

\begin{lstlisting}
public class InvalidVehicleDataException extends Exception {
    public InvalidVehicleDataException(String message) {
        super(message);
    }
}

// Usage
if (vehicleSpeed < 1 || vehicleSpeed > 10) {
    throw new InvalidVehicleDataException(
        "Vehicle speed must be between 1 and 10");
}
\end{lstlisting}

\subsection{File Handling and Logging}
The \texttt{SimulatorLogger} class demonstrates I/O and persistence:

\begin{lstlisting}
public class SimulatorLogger {
    private PrintWriter fileWriter;
    private Queue<String> logBuffer;
    
    public SimulatorLogger(String fileName) throws IOException {
        fileWriter = new PrintWriter(new FileWriter(fileName, true));
        logBuffer = new LinkedList<>();
    }
    
    public void log(String message) {
        String timestampedMsg = "[" + System.currentTimeMillis() 
            + "] " + message;
        fileWriter.println(timestampedMsg);
        logBuffer.add(timestampedMsg);
    }
    
    public void flush() {
        fileWriter.flush();
    }
}
\end{lstlisting}

\section{System Design}
\subsection{Architecture Overview}
The system follows a layered, event-driven architecture:
\begin{itemize}
    \item \textbf{Presentation Layer:} Swing GUI renders roads, vehicles, lights, and controls
    \item \textbf{Simulation Layer:} Timer-driven update loops handle entity movements
    \item \textbf{Control Layer:} \texttt{TrafficController} manages signal logic and lane transitions
    \item \textbf{Domain Layer:} \texttt{Vehicle} hierarchy with type-specific behaviors
    \item \textbf{Utility Layer:} Logging, exceptions, and helper classes
\end{itemize}

\subsection{Package Structure}
\begin{lstlisting}
com/smartcity/traffic/
  Main.java                          # Application entry point
  TrafficSimulatorGUI.java          # 1335 lines - main UI and rendering
  TrafficController.java            # Signal logic and cycling
  TrafficLight.java                 # Signal state machine
  Vehicle.java                      # Abstract vehicle base class
  PrivateCar.java                   # Private car implementation
  PublicBus.java                    # Public bus implementation
  Motorcycle.java                   # Motorcycle implementation
  Ambulance.java                    # Emergency vehicle implementation
  EmergencyVehicle.java             # Interface for emergency behavior
  SimulatorLogger.java              # File logging with buffering
  util/
    InvalidVehicleDataException.java # Custom exception class
\end{lstlisting}

\newpage
\subsection{Block Diagram / Architecture}
The system follows a layered, event-driven architecture:
\begin{itemize}
    \item Presentation Layer: Swing GUI renders roads, vehicles, traffic lights, controls, and statistics.
    \item Simulation Layer: timer-driven update loops handle vehicle movement and signal updates.
    \item Control Layer: \texttt{TrafficController} executes cycle management, smart redistribution, and emergency override.
    \item Domain Layer: \texttt{Vehicle} hierarchy models behavior differences by vehicle type.
    \item Utility Layer: logging and custom exception classes support robustness and observability.
\end{itemize}

(Add diagram below)
\begin{figure}[h]
    \centering
    %\includegraphics[width=0.8\textwidth]{diagram.png}
    \framebox(300,150){Place Block Diagram Here} % Placeholder box
    \caption{System Architecture}
\end{figure}

\subsection{Modules}
\begin{itemize}
    \item Traffic Control Module (\texttt{TrafficController}, \texttt{TrafficLight}):
    controls cycle sequencing, yellow transitions, lane switching, time-bank usage, and emergency preemption.
    \item Vehicle Module (\texttt{Vehicle}, \texttt{PrivateCar}, \texttt{PublicBus}, \texttt{Motorcycle}, \texttt{Ambulance}):
    defines lane-based motion rules, speed profiles, and special movement behavior.
    \item User Interface Module (\texttt{TrafficSimulatorGUI}):
    provides spawn controls, status indicators, simulation speed slider, pause/resume, and visualization.
    \item Logging and Utilities Module (\texttt{SimulatorLogger}, \texttt{InvalidVehicleDataException}):
    provides file logging, in-memory log buffering, and custom validation support.
    \item Bootstrap Module (\texttt{Main}):
    initializes logger, sets look-and-feel, and launches the GUI.
\end{itemize}

\section{Implementation Details}
The project is implemented as a desktop Java application with clear separation between simulation logic and presentation logic.

\subsection{Core Traffic Control Algorithm}
The \texttt{TrafficController} maintains the following state machine:
\begin{itemize}
    \item \textbf{Fixed Cycle:} 80 seconds total, 20 seconds per lane (North → East → South → West → repeat)
    \item \textbf{Signal States:} Each lane's light progresses RED → GREEN → YELLOW → RED
    \item \textbf{Smart Mode:} When active, checks if the current lane is empty (0 vehicles). If empty for 3+ seconds, the system saves time and moves to the next lane early, accumulating the difference in a time bank.
    \item \textbf{Time Bank Usage:} When a heavily congested lane (5+ vehicles) becomes active, the controller extends its green phase by withdrawing time from the bank, up to a maximum of +10 seconds.
    \item \textbf{Emergency Override:} When an ambulance is detected by the GUI, \texttt{trafficController.handleEmergencyOverride(laneID)} is called. This immediately sets that lane to GREEN and overrides the normal cycle, giving the ambulance priority.
\end{itemize}

\subsection{Vehicle Movement Model}
Each vehicle implements \texttt{move(String signalState)} differently:
\begin{itemize}
    \item \textbf{RED:} No vehicle of any type moves.
    \item \textbf{GREEN:} All vehicles move forward at their configured speed.
    \item \textbf{YELLOW:} Different types have different crossing probabilities:
    \begin{itemize}
        \item Private Car: 30\% chance to proceed (conservative)
        \item Bus: 40\% chance (slightly more adventurous than car)
        \item Motorcycle: 70\% chance (aggressive)
        \item Ambulance: 100\% chance when emergency is active
    \end{itemize}
    \item \textbf{Collision Avoidance:} Vehicles check the distance to the next vehicle ahead and reduce speed if a collision is imminent.
\end{itemize}

\subsection{GUI and Real-Time Updates}
The \texttt{TrafficSimulatorGUI} uses two concurrent Swing timers:
\begin{itemize}
    \item \textbf{Animation Timer (50ms interval):} Calls \texttt{vehicle.move()} for each vehicle and triggers a panel repaint. This provides smooth animation.
    \item \textbf{Control Timer (1000ms interval):} Advances the signal countdown, checks for lane switching conditions, updates status labels, and manages the overall cycle state.
\end{itemize}
This separation ensures that rendering is smooth while control logic remains deterministic.

\subsection{File Logging and Session Management}
The \texttt{SimulatorLogger} creates timestamped log files:
\begin{itemize}
    \item Each session is logged to a file under \texttt{logs/} directory
    \item Lines are timestamped with millisecond precision
    \item Recent log entries (last 100) are held in memory for quick viewing via the GUI's ``View Logs'' button
    \item Example log entries:
    \begin{itemize}
        \item \texttt{[1712761234567] SIMULATION STARTED - Smart Mode: enabled}
        \item \texttt{[1712761245890] NORTH lane GREEN (20s) - Vehicle count: 5}
        \item \texttt{[1712761301234] Emergency ambulance detected on WEST - overriding signal}
        \item \texttt{[1712761412567] SIMULATION PAUSED}
    \end{itemize}
\end{itemize}

\section{Sample Outputs and Usage Scenarios}

\subsection{Scenario 1: Fixed-Cycle Mode with Standard Traffic}
\begin{enumerate}
    \item User starts the application; all lanes are initially empty (RED lights).
    \item User clicks ``Spawn Vehicles (NORTH)'' three times  → 3 private cars appear in the NORTH lane queue.
    \item Signal controller immediately switches NORTH to GREEN (20s countdown begins).
    \item All three cars move forward and pass through the intersection.
    \item After 20 seconds, NORTH changes to YELLOW (3s countdown), then RED.
    \item EAST lane signal turns GREEN; cycle continues.
    \item \textbf{Expected Output:} Lane counts display [NORTH: 0, EAST: 0, ...], cycle timer shows active lane and countdown.
\end{enumerate}

\subsection{Scenario 2: Smart Mode with Empty Lane}
\begin{enumerate}
    \item User clicks ``Toggle Smart Mode'' → displays ``Smart Mode: ON''.
    \item User spawns 5 vehicles in NORTH and 1 vehicle in EAST.
    \item NORTH is active (GREEN for 20s), EAST is waiting (RED).
    \item After 5 seconds with no new cars in EAST and NORTH empty, smart mode detects the empty EAST lane.
    \item System saves 3 seconds and switches to EAST (Green). Saved time = 3 seconds → time bank shows ``+3s''.
    \item EAST green lasts 20 - 3 = 17 seconds. One vehicle exits.
    \item When NORTH becomes active again later, it receives 20 + 1 = 21 seconds (uses 1s from time bank).
    \item \textbf{Expected Output:} Time bank label shows accumulated or used time; lane switch happens before full cycle.
\end{enumerate}

\subsection{Scenario 3: Emergency Override}
\begin{enumerate}
    \item User spawns standard traffic in all lanes (2-3 cars each).
    \item Currently SOUTH is GREEN, WEST is RED.
    \item User clicks ``Spawn Ambulance (WEST)'' → an ambulance appears in the WEST lane with red color and flashing border.
    \item GUI detects ambulance and calls \texttt{handleEmergencyOverride("WEST")}.
    \item System immediately sets SOUTH to YELLOW → RED, and WEST to GREEN.
    \item Ambulance moves forward at max speed (8 px/frame) and exits the intersection.
    \item Log displays: ``Emergency override activated - WEST lane priority granted''.
    \item ``Emergency Count'' label increments to 1.
    \item After ambulance clears, normal cycle resumes.
    \item \textbf{Expected Output:} Signal state changes instantly; ambulance crosses while other lanes are red; counter increments.
\end{enumerate}

\subsection{Scenario 4: Pause and Resume}
\begin{enumerate}
    \item Simulation is running with active traffic (vehicles moving, lights changing).
    \item User clicks ``Pause Simulation'' → all timers stop, vehicle motion freezes.
    \item Countdown timers display frozen values; no log updates occur.
    \item User clicks ``Resume Simulation'' → timers restart from the paused state (not reset).
    \item Traffic continues as if nothing happened.
    \item \textbf{Expected Output:} Smooth pause/resume with state preservation; no reset of cycle or vehicle positions.
\end{enumerate}

\subsection{Scenario 5: View Logs}
\begin{enumerate}
    \item After a few minutes of simulation, user clicks ``View Logs''.
    \item A dialog window opens showing the most recent 100 timestamped log entries.
    \item Entries show sequence of signal changes, vehicle counts, emergency events, mode toggles, etc.
    \item User can scroll to see older logs from the current session.
    \item Log file is simultaneously persisted to disk.
    \item \textbf{Expected Output:} Dialog displays formatted, scrollable timestamped log records; no performance degradation.
\end{enumerate}

\section{UML Class Diagram}
The diagram below illustrates the complete class hierarchy and relationships:

\begin{figure}[h]
    \centering
    \framebox(400,200){[UML Class Diagram to be inserted here]}
    \caption{Complete UML Class Hierarchy showing Vehicle inheritance, EmergencyVehicle interface, and TrafficController dependencies}
\end{figure}

Key relationships:
\begin{itemize}
    \item \texttt{Vehicle} is abstract; all concrete vehicles inherit from it.
    \item \texttt{Ambulance} implements both \texttt{Vehicle} inheritance AND \texttt{EmergencyVehicle} interface.
    \item \texttt{TrafficController} maintains a \texttt{HashMap<String, TrafficLight>} for lane lookup.
    \item \texttt{TrafficSimulatorGUI} holds references to lists of \texttt{Vehicle} objects and the \texttt{TrafficController}.
    \item \texttt{SimulatorLogger} is instantiated once in \texttt{Main} and used globally for logging.
\end{itemize}

% ==========================================
% Previous SCREENSHOTS section removed, replaced with Sample Outputs above
% ==========================================

\section{Testing}
\begin{table}[h]
    \centering
    \begin{tabularx}{\textwidth}{|l|X|X|l|}
        \hline
        \textbf{Test Case} & \textbf{Input} & \textbf{Expected Output} & \textbf{Result} \\ \hline
        TC-01 & Spawn normal traffic in one lane & Vehicles appear and move according to that lane state & Pass \\ \hline
        TC-02 & Spawn traffic in multiple lanes & Per-lane counts update correctly & Pass \\ \hline
        TC-03 & Wait for cycle transition & Green to yellow to next lane transition occurs correctly & Pass \\ \hline
        TC-04 & Enable smart mode with empty active lane & Early switch occurs and saved time is accumulated & Pass \\ \hline
        TC-05 & Enable smart mode with congested lane & Saved time is reused as green extension for busy lane & Pass \\ \hline
        TC-06 & Spawn ambulance on non-green lane & Emergency override shortens/adjusts current phase & Pass \\ \hline
        TC-07 & Pause simulation & Vehicle motion and controller updates stop temporarily & Pass \\ \hline
        TC-08 & Resume simulation & Updates continue from paused state without reset & Pass \\ \hline
        TC-09 & Open log viewer & Recent timestamped logs appear in GUI dialog & Pass \\ \hline
        TC-10 & Quit application button & Timers stop and application exits cleanly & Pass \\ \hline
    \end{tabularx}
\end{table}

\section{Results and Discussion}
The simulator produced stable and consistent behavior across all defined test cases. Signal transitions followed the expected finite-state progression, vehicle movement respected lane and signal constraints, and emergency prioritization was triggered reliably when ambulances were introduced.

From a control perspective, the smart-mode logic reduced idle green duration for empty lanes and reallocated saved time to lanes with higher load, which improved visual throughput in high-density scenarios. The lane counters, cycle timer, and emergency override count provided transparent runtime observability and helped validate internal controller decisions.

From an academic perspective, the project successfully demonstrates practical application of OOPJ concepts in a complete, interactive system rather than isolated code snippets. The class hierarchy, interface usage, modular controller design, timer-driven event flow, and logging integration collectively satisfy the major learning outcomes expected in a mini-project.

\section{Limitations}
\begin{itemize}
    \item Simplified single-intersection model without city-wide network behavior
    \item Adaptive logic is rule-based; no machine-learning or predictive optimization
    \item Basic persistence using files only; no full database integration
    \item Vehicle behavior does not include accidents, roadblocks, weather, or pedestrian crossings
    \item Simulation physics are abstracted and not calibrated with real sensor data
\end{itemize}

\section{Future Enhancements}
\begin{itemize}
    \item Integrate multiple intersections with coordinated signal control
    \item Add AI/ML-based adaptive signal timing based on live traffic density
    \item Add pedestrian phases, zebra-crossing control, and violation detection
    \item Integrate camera or IoT sensor inputs for realistic lane-density estimation
    \item Introduce complete database support for long-term analytics
    \item Enhance GUI with charts, heatmaps, and real-time performance metrics
    \item Export simulation reports (CSV/PDF) for comparative scenario analysis
\end{itemize}

\section{Conclusion}
This work successfully implemented a Smart City Traffic Simulator using Java and Swing as a complete OOPJ mini-project. The final system integrates domain modeling, controller logic, real-time GUI interaction, emergency handling, and persistent logging in a single coherent architecture.

Technically, the project demonstrates that even a simplified software model can capture important traffic-management behaviors such as adaptive lane switching and emergency preemption. Academically, it reinforces essential OOP practices including abstraction, inheritance, polymorphism, interface-based design, modular decomposition, and exception-oriented validation.

Overall, the project achieves its core objectives and establishes a strong baseline for future expansion toward data-driven, multi-intersection smart-city traffic systems.

\newpage

% ==========================================
% TEAM MEMBER CONTRIBUTION
% ==========================================
\section{Team Member Contribution}

\begin{table}[h]
    \centering
    \begin{tabularx}{\textwidth}{|l|l|X|}
        \hline
        \textbf{Member Name} & \textbf{Roll No} & \textbf{Key Contributions} \\ \hline
        Anupam & [To be filled] & 
        Vehicle hierarchy design (Vehicle base class, PrivateCar, Motorcycle, Ambulance subclasses); 
        Abstract class design patterns; 
        Core domain modeling. \\ \hline
        
        Arjun & [To be filled] & 
        TrafficController logic (cycle management, smart redistribution, time-bank implementation); 
        TrafficLight state machine; 
        Emergency override algorithm. \\ \hline
        
        Abhinav & [To be filled] & 
        TrafficSimulatorGUI (1335 lines – complete Swing implementation); 
        Real-time rendering and vehicle visualization; 
        Control panel and event listeners; 
        Pause/resume mechanics. \\ \hline
    \end{tabularx}
    \caption{Team Member Contributions}
\end{table}

\subsection{Collaborative Efforts}
\begin{itemize}
    \item \textbf{Logging Module:} All team members contributed to \texttt{SimulatorLogger} design and usage patterns for observability.
    \item \textbf{Testing and Validation:} Comprehensive test suite (10 test cases) defined and executed collaboratively.
    \item \textbf{Documentation:} This report was jointly authored with contributions from all members.
    \item \textbf{Demonstrations:} Code was tested and demonstrated live during integration sessions.
\end{itemize}

\newpage

\section{Bibliography and References}

\begin{enumerate}
    \item[{[1]}] Oracle Corporation. \textit{Java Platform, Standard Edition 17 Documentation}. 2025.
    Available at: \url{https://docs.oracle.com/en/java/javase/17/}
    
    \item[{[2]}] Oracle Corporation. \textit{Java Swing Tutorial}. 2025.
    Available at: \url{https://docs.oracle.com/javase/tutorial/uiswing/}
    
    \item[{[3]}] Oracle Corporation. \textit{Java Collections Framework}. 2025.
    Available at: \url{https://docs.oracle.com/javase/8/docs/technotes/guides/collections/}
    
    \item[{[4]}] Erich Gamma, Richard Helm, Ralph Johnson, John Vlissides. 
    \textit{Design Patterns: Elements of Reusable Object-Oriented Software}. 
    Addison-Wesley Professional, 1st Edition, 1994.
    
    \item[{[5]}] Joshua Bloch. 
    \textit{Effective Java: Programming Language Guide}. 
    Addison-Wesley, 3rd Edition, 2018.
    
    \item[{[6]}] Horstmann, Cay S., and Gary Cornell. 
    \textit{Core Java: Fundamentals}. 
    Prentice Hall, 12th Edition, 2021.
    
    \item[{[7]}] Sayed Z. Nekoo, et al. 
    \textit{Adaptive Traffic Signal Control: A Comprehensive Review}. 
    IEEE Transactions on Intelligent Transportation Systems, Vol. 22, No. 8, 2021.
    
    \item[{[8]}] Project Repository and Source Code -- GitHub (with CI/CD workflows and documentation).
    
    \item[{[9]}] OOPJ Lab Mini-Project Guidelines and Academic Rubric.
\end{enumerate}

\end{document}